%! Mode:: "TeX:UTF-8"
%!TEX encoding = UTF-8 Unicode
%!TEX TS-program = xelatex
\documentclass{ctexart}
\usepackage{bbm}
\usepackage{tikz}
\usepackage{amsmath,amssymb,amsthm,color,mathrsfs}
\usepackage{hyperref}
\usepackage{cleveref}
\usepackage{enumitem,anysize}%
\usepackage{expl3}
\marginsize{1in}{1in}{1in}{1in}%
\theoremstyle{remark}
\newtheorem*{claim*}{Claim}
\newtheorem{lemma}{Lemma}

\DeclareMathOperator{\ot}{ordertype}
\DeclareMathOperator{\dom}{dom}
\DeclareMathOperator{\ran}{ran}
\DeclareMathOperator{\field}{field}
\DeclareMathOperator{\Ord}{Ord}

\newcommand{\cc}{\mathfrak{c}}

\def\email#1{{\texttt{#1}}}
\newcommand{\isep}[1][0pt]{\addtolength{\itemsep}{#1}}
\newcounter{problem}
\renewcommand{\theproblem}{\Roman{problem}}
\newenvironment{problem}{\refstepcounter{problem}\noindent\color{blue}\begin{tikzpicture}[baseline]%
\node at (-0.02em,0.3em) {$\mathbb{P}$};%
\node[scale=0.7] at (0.2em,-0.0em) {R};%
\node[scale=0.7] at (0.6em,0.4em) {O};%
\node[scale=0.8] at (1.05em,0.25em) {B};%
\node at (1.55em,0.3em) {L};%
\node[scale=0.7] at (1.75em,0.45em) {E};%
\node at (2.35em,0.3em) {M};%
\end{tikzpicture}\theproblem}{}
\crefname{problem}{Problem}{Problem}
%\renewcommand\theprob{{\Roman{problem}}}
\newcommand\mysolution{\begin{tikzpicture}[baseline]%
\node at (-0.04em,0.3em) {$\mathbb{S}$};%
\node[scale=0.7] at (0.35em,0.4em) {O};%
\node at (0.7em,0.3em) {\textit{L}};%
\node[scale=0.7] at (0.95em,0.4em) {U};%
\node[scale=1.1] at (1.19em,0.32em){T};%
\node[scale=0.85] at (1.4em,0.24em){I};%
\node at (1.9em,0.32em){$\mathcal{O}$};%
\node[scale=0.75] at (2.3em,0.21em){\texttt{N}};%
\end{tikzpicture}}
\newenvironment{solution}{\begin{proof}[\mysolution]}{\end{proof}}

%%%%%%%%
\newcommand{\calL}{\mathcal{L}}

\newcommand\<{\langle}
\renewcommand\>{\rangle}
\newcommand\eneg{\mathcal{E}_{\neg}}
\newcommand\eto{\mathcal{E}_{\to}}
\newcommand\N{\mathbb{N}}
\newcommand\subini{\subsetneqq_{init}}
\def\to{\rightarrow}
\newcommand{\calA}{\mathcal{A}}
\newcommand{\xor}{\vee}
\newcommand{\bor}{\bigvee}
\newcommand{\band}{\bigwedge}
\newcommand{\xand}{\wedge}
\newcommand{\minus}{\mathbin{\backslash}}
\newcommand{\calF}{\mathcal{F}}
\newcommand\mi[1]{\mathscr{P}(#1)}
\newcommand{\calB}{\mathcal{B}}
\newcommand{\heq}{\mathop{\hat{=}}}
\newcommand{\hneq}{\mathop{\hat{\neq}}}
\newcommand{\calR}{\mathcal{R}}
\newcommand{\hle}{\mathop{\hat{<}}}
\newcommand{\R}{\mathbb{R}}
\newcommand{\calM}{\mathcal{M}}
\newcommand{\calN}{\mathcal{N}}
\newcommand{\frA}{\mathfrak{A}}
\newcommand{\card}{\mathrm{card}}
\newcommand{\frM}{\mathfrak{M}}
\newcommand{\Th}{\mathrm{Th}}
\newcommand{\ecd}{\mathrm{EC}_\Delta}
\newcommand{\Q}{\mathbb{Q}}
\newcommand{\Z}{\mathbb{Z}}
\newcommand{\oto}{\leftrightarrow}
\newcommand{\hin}{\hat{\in}}
\newcommand{\fun}[2]{{}^{#1}#2}
\newcommand{\A}{\mathbbm{A}}








\begin{document}
\title{$\mathbb{ALGEBRAIC\ GEOMETRY}$}
\author{白永乐\\ %% 姓名
SID: 202011150087\\ %% 学号
\email{202011150087@mail.bnu.edu.cn}} %% 电邮
\maketitle
\begin{problem}
    $P$ is an ideal of a unitary commutative ring $A$, then $P$ is prime ideal of $A\iff A/P$ is integral domain.
\end{problem}
\begin{solution}
    $\Rightarrow $:

    Since $A$ is a unitary commutative ring, so $A/P$ is unitary commutative ring, too. So we only need to prove $[ab]=[0]\Rightarrow [a]=[0]\xor [b]=[0]$. 
    Obviously $[ab]=0\iff ab\in P\iff a\in P\xor b\in P\iff [a]=[0]\xor [b]=[0]$. 

    $\Leftarrow $:

    As the same, $ab\in P\iff [ab]=[0]\Rightarrow [a]=[0]\xor [b]=[0]\iff a\in P\xor b\in P$, so $P$ is prime ideal. 
\end{solution}

\begin{problem}
    $M$ is an ideal of a unitary commutative ring $A$, then $M$ is maximal ideal of $A \iff A/M$ is a field.   
\end{problem}
\begin{solution}
    $\Rightarrow $:

    Consider $[a]\in A/M\minus [0]$, we will prove it has a reverse.  Consider $N:=\{xm+ya:x,y\in A,m\in M\}$ is the minimum ideal of $A$ contains $M$ and $a$. 
    Since $[a]\neq [0]$ we know $a\notin M$, so $M\subsetneqq N$. 
    Noting $M$ is maximal, so $N=A$. That means $\exists x,y\in A,m\in M,xm+ya=1$. So $[xm+ya]=[1]$. Since $[xm]=[0]$ we get $[y][a]=1$, i.e., $[y]=[a]^{-1}$. 

    $\Leftarrow $:

    Consider $a\in A\minus M, N:=\{xp+ya:x,y\in A,p\in P\}$, we will prove $N=A$, which means $M$ is maximal. 
    Since $A/M$ is field, $\exists y\in A, [y]=[a]^{-1}$. Thats means $ay-1\in M\subset N$. Noting $ay\in N$, so $1\in N$, thus $N=A$. 
\end{solution}

\begin{problem}
    A ring $A$ is noetherian, $I\subset A$ is an ideal of $A$, then $A/I$ is noetherian. 
\end{problem}
\begin{solution}
    Consider an ideal $J\subset A/I$, let $M:=\{x\in A:[x]\in J\}$. Then $\forall a\in A,x\in M,[ax]=[a][x]\in J$, so $ax\in M$. $\forall a,b\in M,[a-b]=[a]-[b]\in J$, so $a-b\in M$. So $M$ is an ideal of $A$. 
    Since $A$ is noetherian, we can assume $M=(f_i,i=1,2,\cdots n)$. Now we will prove $J=([f_i],i=1,2,\cdots n)$. 
    Consider $[f]\in J$, from definition of $M$ we know $f\in M$, so $f=\sum_{i=1}^na_if_i,a_i\in A$, thus $[f]=\left[\sum_{i=1}^na_if_i\right]=\sum_{i=1}^n[a_i][f_i]$. So $J=([f_i],i=1,2,\cdots n)$. 
\end{solution}

\begin{problem}
    \label{pro:4}
    $K$ is a field, $A=K[x_1,x_2,\cdots x_n],\A^n_K=K^n$. For ideal $I$ of $A$ let $V(I):=\{p\in\A_K^n:f(p)=0,\forall f\in I\}$. Then $V(I_1)\cup V(I_2)=V(I_1\cap I_2)=V(I_1I_2)$
\end{problem}
\begin{solution}
    \begin{lemma}
        \label{lem:1}
        $I\subset J\Rightarrow V(I)\supset V(J)$.
    \end{lemma}
    \begin{proof}
        Consider $p\in V(J)$, we get $\forall f\in J,f(p)=0$. Since $I\subset J$, so $\forall f\in I,f(p)=0$, i.e., $f\in V(I)$. 
    \end{proof}
    From Lemma \Cref{lem:1} we know $V(I_1\cap I_2)\subset V(I_1I_2)$, so we only need to prove $V(I_1)\cup V(I_2)\subset V(I_1\cap I_2),V(I_1I_2)\subset V(I_1)\cup V(I_2)$. 
    \begin{itemize}
        \item $V(I_1)\cup V(I_2)\subset V(I_1\cap I_2)$:
        From \Cref{lem:1} It's obvious. 
        \item $V(I_1I_2)\subset V(I_1)\cup V(I_2)$:
        Consider $p\in V(I_1I_2)$. 
        If $p\notin V(I_1)\cup V(I_2)$, then $\exists f_1\in I_1,f_2\in I_2, f_1(p)\neq 0,f_2(p)\neq 0$. Now consider $f=f_1f_2\in I_1I_2$, we get $f(p)=f_1(p)f_2(p)\neq 0$, so $p\notin V(I_1I_2)$, it's a contradiction. 
    \end{itemize}
\end{solution}

\begin{problem}
    $K$ is an infinite field, then $\A_k^n$ is not Hausdorff. 
\end{problem}
\begin{solution}
    \begin{lemma}
        \label{lem:2}
        $K$ is a infinite field, $f\in K[x_1,x_2,\cdots x_n]\minus \{0\}$, then exists $p\in \A_K^n,f(p)\neq 0$. 
    \end{lemma}
    \begin{proof}
        Use MI. When $n=0,K[x_1,x_2,\cdots x_n]=K$, so it's obvious. Assume for $n=k$ it's right, when goes to $k+1$:

        Consider $h\in K(x_1,x_2,\cdots x_k)[x_{k+1}],h(x_{k+1}):=f(x_1,x_2,\cdots x_k,x_{k+1})$ is a non-zero polynomial so it has finite root. So exists $a\in K,h(a)\neq 0$. So $g:=f(x_1,x_2,\cdots x_k,a)\in K[x_1,x_2,\cdots x_k]\neq 0$. 
        By induction hypothesis we get $\exists b\in \A_K^k,g(b)\neq 0$. Let $p:=(b,a)\in \A_K^{k+1}$, then $f(p)\neq 0$. 
    \end{proof}
    In fact, it's not only not Hausdorff, it's kind of ''absolutly not Hausdorff'' because every pair of point can't be seperated. 
    Consider two point $p\neq q,p,q\in \A_K^n$. 
    Assume $p=(p_1,p_2,\cdots p_n),q=(q_1,q_2,\cdots q_n),p_1\neq q_1$. Assume two open set $V(I_1)^c,V(I_2)^c$ can seperate $p,q$, then $V(I_1)\cup V(I_2)=\A_K^n$. 
    From \Cref{pro:4} we know $V(I_1I_2)=\A_K^n$. So $\forall f\in V(I_1I_2),\forall p\in \A_K^n,f(p)=0$. 
    Then from Lemma\Cref{lem:2} we can get $f=0$. So $I_1I_2=\{0\}$. 
    Since $p\notin V(I_1),q\notin V(I_2)$, we know $I_1,I_2\neq \{0\}$. So $\exists f_1\in I_1\neq 0,\exists f_2\in I_2\neq 0$, and thus $f=f_1f_2\in I_1I_2\neq 0$, contradiction! 
    So these is not a pair of points can be seperated by two disjoint open set. 
\end{solution}

\end{document}
